% Transcription — fonds Alexandre Grothendieck, shelfmark 15, batch 11 (pages 201-220).
% Established from the facsimile archives/batches/15/batch-11.pdf (a mirror of
% https://grothendieck.umontpellier.fr/15.pdf), per the skill
% transcribe-grothendieck. The batch file carries no cover sheet: PDF page N
% is page 200+N of the folder (the Montpellier footer prints « 201 » ... « 220 »),
% and the archivists' pencil numbers bottom-left agree wherever checked (202,
% 206, 208, 212, 214).
%
% Pass: Opus 5.5 (claude-opus-5-5), 2026-09-23 — first pass, unchecked against the pages by a human.
%
% Ten of the twenty pages are transcribed: the even pages 202-220. Absent:
% the odd pages 201-219, leaves of unrelated typescripts — Bourbaki drafts
% (« Tribu n° 59 »: « Algèbres de Lie semi-simples (3), (4) », « Systèmes de
% racines (2), (3) », « Systèmes de racines. Introduction (Rédaction
% Borel) »; a draft « n° 403 » on semi-simple Lie algebras; a draft « n° 413 »
% on Lie groups, pp. 12-19), many scanned upside down, nothing in his hand.
% They are not the folder's own « tapuscrit ».
%
% What the batch is. Ten handwritten leaves, not paginated by him, fast and
% often hard, continuing the fibre functors of motives at a point x
% (ξ_{x,Z}, ξ_{x,R}, ξ_{x,∞}, ξ_{x,Q}, ξ_{x,ℓ}) begun in batch 10:
%   202    Besides the integral structure on the Hodge-De Rham functor at a
%          complex point, the decomposition into types (compatible with the
%          canonical filtration), determined via the filtration by the
%          semi-linear involution τ' on ξ_{x,∞} = ξ_{x,R} ⊗_R C; for x real a
%          second semi-linear involution τ'' from ξ_{x,∞} ≃ ξ_{x,∞}|R ⊗_R C,
%          and τ = τ'τ'' = τ''τ' on (ξ_{x,Z})_C (τ' « transcendante », τ''
%          « algébrique »). The local situation at a real point is known from
%          τ on ξ_{x,Z}, hence on ξ_{x,ℓ}: « or c'est là l'élément de π_1(x),
%          Frobenius en la place ... réelle ! ».
%   204    π_1(x) does not act trivially on ξ_{x,Z}, ξ_{x,R}, but f_∞ (the
%          Frobenius at infinity) does, extended linearly to τ = f_∞. N.B.:
%          the Hodge conjecture read as: the part of level ≤ ρ of ξ_{x,Q}(F)
%          is ξ_{x,Q}(F) ∩ ξ_{x,∞}(F)^{(ρ)} ∩ f_∞ ξ_{x,∞}(F)^{(ρ)}; in
%          particular these are motivic (coming from a sub-motive F^{(ρ)}).
%   206    Varying x: he thinks the ξ_{x,Z}, ξ_{x,R} at distinct places are
%          non-isomorphic (circled margin « faux, cf. groupoïde fondamental
%          sur Z »); the ξ_{x,C} = ξ_{x,∞} are isomorphic, not canonically.
%          N.B. for p > 0, T_α(R) = ∅ (?). Then « Pt motivique de X à valeurs
%          dans un corps (ou un anneau) k »: the objects of T(k); the points
%          form a fibred category, a stack for fpqc; isomorphisms form a
%          torsor (the sentence runs onto 208).
%   208    Points ℓ-adic (values Z_ℓ, Q_ℓ), transcendental (values R or Z,
%          « structure supplémentaire », matching the points of X in C mod
%          conjugation), Hodge-De Rham (values in any field), Witt (W(k), k of
%          char. p).
%   210    « Questions »: points of X in k of char. p > 0 vs motivic points in
%          W(k) with (F, V) (« foncteurs des groupes formels ») → « pts
%          motiviques généralisés »; in char. 0, motivic points with filtered
%          structure; complex points mod conjugation ⟷ motivic points in R,
%          in (C, f_∞), in Z; « Pts banaux » in Z_ℓ (ℓ prime to the residue
%          characteristics) or Q_ℓ; finite fields F_q ⟷ W(F_q) with (F, V).
%   212    Heading « Groupe fondamental motivique discret »: ξ ∈ T(Z),
%          π_1(X,ξ) = Aut(ξ) = G_ξ(Z); N.B. by Borel-(Harish-Chandra?) only
%          finitely many ξ ∈ T_α(Z), hence finitely many discrete fundamental
%          groups. For ξ ∈ T_{(α)}(k), G(X,ξ) = Aut(ξ) (underlined), π_1(X,ξ)
%          ≃ G(X,ξ)(k); comparison with π_1(X^an, ξ); X over C: π_1(X,ξ) with
%          a Hodge structure.
%   214    (from the Hodge and Tate conjectures) « Caractérisation de l'image
%          ?? »; same over k ⊂ C algebraically closed, k = Q̄; « Analyser
%          l'exemple de Serre de π_1 non isomorphes »; base field R: a real
%          point gives ξ_a ∈ T_α(Z) and an involution compatible with complex
%          conjugation on G(X,ξ); the functor M ↦ (ξ(M_C), Hodge structure +
%          real structure by f_∞) fully faithful. Pb: how the Hodge
%          structure varies with the complex point.
%   216    Link with moduli varieties of rigidified polarised motives: a
%          morphism X → M is determined by its value at one point plus π_1(X)
%          → π_1(M). Pointed motivic schemes, morphisms G(X,ξ) → G(Y,η), hence
%          π_1(X,ξ) → π_1(Y,η).
%   218    A palimpsest: X of finite type over k, v, w two archimedean places,
%          fibre functors ξ^Z_v, ξ^Z_w; questions a) isomorphic? b) on the
%          finite-type subcategories M_α(X)? Reduction to k algebraically
%          closed in k(X); 1°) X defined over k_0 (algebraic closure of Q in
%          k): a) ξ^Z_v | M_α(X) ≃ ξ^Z_w | M_α(X) for all α, via X_0 over A_0,
%          k = E_0 and the transcendental fundamental group.
%   220    k countable: k a filtered union of k_0-algebras A_i of finite type,
%          transitions smooth with geometrically connected fibres, π^1(x_j,
%          y_j) → π^1(x_i, y_i) surjective, lim ≠ ∅, cqfd; k uncountable
%          reduced to E. 2) E ≠ k_0: « je pense » some α with ξ^Z_v|M_α(X)
%          ≄ ξ^Z_w|M_α(X); one may suppose X = Spec k, k a finite extension of
%          Q. The rest of the page is cancelled. Batch 12 opens (page 222)
%          with « Il s'agirait de construire X projectif lisse sur k ... »,
%          the continuation presumably.
%
% Notation fixed for the batch, following batch 10. ξ with indices as he
% writes them (ξ_{x,Z}, ξ_{x,R}, ξ_{x,C}, ξ_{x,∞}, ξ_{x,Q}, ξ_{x,ℓ}, ξ^Z_v,
% ξ^Z_w); f_∞ the Frobenius at infinity; τ, τ', τ''; T(k), T_α(k), T_{(α)}(k)
% the points; \mathcal{M}, \mathcal{M}_\alpha(X) for his script M;
% \underline{\mathrm{Aut}} for his underlined Aut. « ½ linéaire »
% (semi-linear), « pt », « pts », « val. », « alg. », « conj. », « ens. »,
% « cqfd » kept as he writes them. Plain square brackets in the text are his
% own brackets or words he added in the interline.
%
% Legibility. The formulas carry; the connective prose of 206, 208 and
% 214 is partly \ill{}; 218 and 220 are heavily reworked (lines struck,
% loops of insertion) and are transcribed only in their surviving layer,
% with the cancelled passages summarised in notes. Nothing on these leaves
% is dated.
%
% The dating is the inventory's own for the folder, copied verbatim. Its
% group, [10-18] « Théorie des motifs », is dated [à partir de 1958-à partir
% de 1982].
\documentclass[11pt,a4paper]{article}
% Le préambule commun à toutes les transcriptions du fonds Grothendieck.
%
% Il tient en une idée : **l'incertitude se compose, elle ne se résout pas.**
% Une transcription qui lisse un mot douteux a détruit la seule chose pour
% laquelle on la faisait — un lecteur doit pouvoir savoir, à chaque endroit,
% ce qui était sur le feuillet et ce qui a été ajouté. Les six macros qui
% suivent sont donc l'appareil critique, et non de la décoration.
%
% Les mêmes commandes sont reconnues par `scripts/render.mjs`, qui produit la
% vue de lecture du site. Une macro ajoutée ici doit l'être aussi là-bas,
% faute de quoi le rendu échouera bruyamment — ce qui est voulu : un rendu
% silencieusement faux est pire qu'un rendu absent.

\ProvidesPackage{grothendieck}[2026/08/08 Transcriptions du fonds Grothendieck]

\usepackage{amsmath,amssymb,amsthm}
\usepackage{mathrsfs}
\usepackage{stmaryrd}
% Les diagrammes commutatifs sont partout dans ce fonds : c'est la langue
% même de la géométrie algébrique de Grothendieck, et une transcription qui
% les rendrait en prose ne transcrirait rien.
\usepackage{tikz-cd}
% Français par défaut — c'est la langue des feuillets — mais l'anglais est
% chargé aussi : la traduction et le résumé sont des documents anglais, et la
% typographie française (espaces avant les deux-points) y serait une faute.
% Ces documents basculent par \selectlanguage{english} en tête de corps.
\usepackage[english,french]{babel}
\usepackage{fontspec}
\usepackage{geometry}
\geometry{a4paper,margin=25mm,marginparwidth=28mm}
\usepackage{marginnote}
\usepackage{xcolor}
% ulem, not soul. `soul`'s \st and \ul are built on a character-by-character
% scan that XeLaTeX's font machinery breaks: the document fails with an
% undefined control sequence rather than degrading. ulem does the same two jobs
% and is Unicode-engine safe. `normalem` stops it hijacking \emph, which the
% transcriptions use for Grothendieck's own emphasis.
\usepackage[normalem]{ulem}
\usepackage{hyperref}
\hypersetup{colorlinks=true,linkcolor=black,urlcolor=black,pdfborder={0 0 0}}

% --- Métadonnées du lot -----------------------------------------------------
% Reprises telles quelles par le script de rendu, qui les affiche en tête de
% la vue de lecture. Elles fixent ce que le fichier prétend couvrir : sans
% elles, une transcription flotte sans point d'attache dans un fonds de seize
% mille feuillets.

\newcommand{\folder}[1]{\gdef\grothFolder{#1}}
\newcommand{\batch}[1]{\gdef\grothBatch{#1}}
\newcommand{\leaves}[2]{\gdef\grothFirst{#1}\gdef\grothLast{#2}}
\newcommand{\foldertitle}[1]{\gdef\grothFoldertitle{#1}}
\gdef\grothFolder{?}\gdef\grothBatch{?}\gdef\grothFirst{?}\gdef\grothLast{?}\gdef\grothFoldertitle{}

% --- Le résumé --------------------------------------------------------------
%
% Il ouvre la lecture modernisée, sous le titre « Résumé », et il est composé
% en petit. Ce n'est pas un ornement : c'est le seul passage du document qui
% ne soit pas des mathématiques, et un lecteur doit pouvoir le reconnaître
% avant de l'avoir lu — puis le sauter s'il connaît déjà le sujet, ou s'y
% arrêter s'il ne le connaît pas. Le filet qui le suit marque la reprise du
% corps.

\newenvironment{resume}
  {\small\setlength{\parskip}{0.45em}}
  {\par\medskip\hrule\medskip}

% --- Mots-clés ---------------------------------------------------------------
%
% En anglais, en clôture du résumé de la lecture modernisée : le vocabulaire
% moderne sous lequel chercher ce que les feuillets construisent. Le manifeste
% du site (`npm run manifest`) les extrait pour étiqueter la cote entière —
% cette ligne est la source unique des tags, il n'y a pas de fichier à côté.
\newcommand{\keywords}[1]{\par\smallskip\noindent
  {\footnotesize\textsc{Keywords} --- \textit{#1}}\par}

% --- L'appareil critique ----------------------------------------------------

% Le numéro de feuillet, en marge. C'est l'ancre qui permet de régler un
% désaccord sur une formule en regardant *un* feuillet plutôt qu'une chemise
% de six cents. Le site s'en sert aussi pour tourner les pages du fac-similé
% quand on fait défiler la transcription.
\newcommand{\leaf}[1]{%
  \marginnote{\footnotesize\color{gray!70}\textsf{#1}}%
  \label{leaf:#1}\ignorespaces}

% Illisible : on ne devine pas. Un mot inventé qui se lit comme les autres est
% le pire résultat possible.
\newcommand{\ill}{\textcolor{red!60!black}{[\,\ldots\,]}}

% Lecture douteuse : le mot est donné, et signalé comme incertain.
\newcommand{\uncertain}[1]{\uline{#1}}

% Ajout de l'éditeur — une lettre manquante, un mot sous-entendu.
\newcommand{\add}[1]{\textcolor{blue!50!black}{[#1]}}

% Biffé par Grothendieck lui-même. Ce qu'il a rayé fait partie du document :
% une rature dit souvent où la pensée a changé de direction.
\newcommand{\struck}[1]{\sout{#1}}

% Note du transcripteur, sur l'état matériel du feuillet.
\newcommand{\note}[1]{\par\smallskip{\small\color{gray!60!black}\textit{[#1]}}\par\smallskip}

% Note marginale de Grothendieck, distincte du corps du texte.
\newcommand{\marginal}[1]{\par\smallskip\noindent\hspace{1em}%
  {\small\color{gray!60!black}#1}\par\smallskip}

% --- Titre ------------------------------------------------------------------

\AtBeginDocument{%
  \begin{center}
    {\large\bfseries Fonds Alexandre Grothendieck --- cote n\textsuperscript{o}~\grothFolder}\\[2pt]
    {\itshape\grothFoldertitle}\\[4pt]
    {\small feuillets \grothFirst--\grothLast\ (lot \grothBatch)}\\[2pt]
    {\footnotesize\color{gray!60!black}Transcription établie d'après le fac-similé numérisé
     par l'Université de Montpellier.\\
     Les lectures incertaines sont soulignées, les illisibles notées [\,\ldots\,],
     les ajouts entre crochets.}
  \end{center}
  \vspace{1em}}

\endinput


\folder{15}
\batch{11}
\pages{201}{220}
\dating{1965-[vers 1977]}
\watermark{Édition de démonstration}
\foldertitle{Théorie arithmétique. Théorie de Galois des motifs (1965) : notes manuscites (s.d.), tapuscrit (s.d.).}

\begin{document}

\page{202}
En plus de la \emph{structure entière} sur le \struck{\ill{}} [foncteur] de
Hodge-De Rham, relatif à un pt complexe, il faut signaler la
\emph{décomposition en types}, i.e.\ \uncertain{graduation} de nature
transcendante (elle est compatible avec la filtration canonique).
\uncertain{Une telle} \ill{} [est déterminée via filtration] par cette
dernière \struck{\ill{}} [par l'opération] \emph{involution} [$\tau$
½ linéaire] $\tau'$ dans $\xi_{x,\infty}$, provenant de l'écriture
\[
  \xi_{x,\infty} = \xi_{x,\mathbb{R}} \otimes_{\mathbb{R}} \mathbb{C}
  \quad\text{où}\quad
  \xi_{x,\mathbb{R}} \simeq \xi_{x,\mathbb{Z}} \otimes_{\mathbb{Z}}
  \mathbb{R} .
\]
\note{« par l'opération » est écrit au-dessus de la ligne, en face d'un mot
biffé ; « $\tau$ ½ linéaire » est ajouté au-dessus de « involution »}

Si $k(x)$ est réel, on a une autre \uncertain{application} ½ linéaire
involutive $\tau''$, provenant de
\[
  \xi_{x,\infty} \simeq \xi_{x,\infty\,\mathbb{R}} \otimes_{\mathbb{R}}
  \mathbb{C}
\]
et $\tau$ sur $(\xi_{x,\mathbb{Z}})_{\mathbb{C}}$ n'est autre que
$\tau'\tau''$
\[
  \tau = \tau'\tau'' = \tau''\tau'
\]
\note{sous $\tau'$ il écrit « ½ linéaire transcendante », sous $\tau''$
« ½ linéaire algébrique (dans le cas réel) »}
\marginal{à vérifier}

La situation locale [(dans le cas réel)] est connue par la connaissance
[Hodge-De Rham réel] du $\xi_{x,\infty}^{\tau''} = (\xi_{x,\infty}
|\mathbb{R})$, muni de $\bar{\tau} = \tau'$ involution \ill{},
[qui est $\mathbb{R}$-linéaire sur le dit espace], ou mieux par celle de
$\tau = \tau''$ opérant sur $\xi_{x,\infty}^{\tau'} \simeq
\xi_{x,\mathbb{R}}$, i.e.\ \ill{} de $\tau$ sur $\xi_{x,\mathbb{Z}}$, dont on
déduit \ill{} [l'opération de]
\[
  \boxed{\tau \text{ sur } \xi_{x,\ell}}
\]
[or c'est là l'élément de $\pi_1(x)$, \emph{Frobenius} en la place
\uncertain{à l'infini réelle} !]
\marginal{à vérifier}
\note{un trait vertical dans la marge gauche marque ces deux dernières
lignes ; le « à vérifier » de la marge est répété à leur hauteur}

\page{204}
Ainsi, $\pi_1(x)$ n'opère trivialement ni sur $\xi_{x,\mathbb{Z}}$, ni
\uncertain{a fortiori} sur $\xi_{x,\mathbb{R}}$, mais l'élément $f_\infty$ de
$\pi_1(x)$ (Frobenius infini) opère sur $\xi_{x,\mathbb{Z}}$ et a fortiori sur
$\xi_{x,\mathbb{R}}$ ; on le prolonge \emph{linéairement} en $\tau =
f_\infty$ …

\emph{N.B.} La \emph{conjecture de Hodge} interprète la filtration
\emph{induite} sur la \uncertain{coh.} rationnelle transcendante
$\xi_{x,\mathbb{Z}} \otimes_{\mathbb{Z}} \mathbb{Q}$ [$=
\xi_{x,\mathbb{Q}}$] par la filtration arithmétique (ou motivique, par le
niveau), en termes de la position \ill{} : la graduation [i.e.\ un terme de
la filtration de Hodge --- \ill{} multiplicatif --- et du $f_\infty$]. Elle
s'exprime en disant que

\struck{Le niveau $\leq \rho$ $\Longleftrightarrow$}
\emph{la partie de niveau} $\leq \rho$ \emph{dans}
$\xi_{x,\mathbb{Q}}(F)$ \emph{n'est autre que}
\[
  \xi_{x,\mathbb{Q}}(F) \cap \xi_{x,\infty}(F)^{(\rho)} \cap f_\infty\,
  \xi_{x,\infty}(F)^{(\rho)} .
\]

En particulier, les derniers \uncertain{termes} sont \uncertain{intégraux}
[?] de $\xi_{x,\mathbb{Q}}(F)$ sont \emph{motiviques} (provenant d'un
sous-motif $F^{(\rho)}$ de $F$), dès que $\xi_{x,\infty}(F)^{(\rho)}$ et
$f_\infty\,\xi_{x,\infty}(F)^{(\rho)}$ sont motiviques …
\struck{\ill{} $\rho$, de dimension \ill{} « niveau »}
\note{la phrase « En particulier … » est peu sûre dans sa construction ; deux
lignes biffées la suivent, dont on ne lit que « $\rho$ » et « niveau », entre
guillemets}

\page{206}
\marginal{faux, cf.\ \uncertain{groupoïde} fondamental sur $\mathbb{Z}$}
\note{note oblique de sa main dans la marge gauche, entourée d'un ovale qui
touche les lignes 3 à 7}

\emph{Faisons varier $x$} pt complexe (\uncertain{ou plus} \ill{}
archimédien). Je pense que \emph{les} $\xi_{x,\mathbb{Z}}$ \emph{et les} $\xi_{x,\mathbb{R}}$ [$\simeq \xi_{x,\mathbb{C}}$ muni de $\tau$]
\emph{qui correspondent à des places distinctes} \emph{sont} \emph{non
isomorphes}. Quant aux $\xi_{x,\mathbb{C}} = \xi_{x,\infty}$, eux sont,
\struck{isomorphes} [\ill{} \uncertain{il est}] \struck{\ill{} comme
\uncertain{nous avons} dit, il faut} \uncertain{alg.\ clos} ils sont
isomorphes entre eux, mais non canoniquement …

\emph{NB} Si on a $p > 0$, je pense que $T_\alpha(\mathbb{R}) = \emptyset$
en général, ce qui explique qu'il n'y ait pas \ill{} \ill{} de
\uncertain{plusieurs} des « composantes connexes » \ill{} « Ramification en
$\mathbb{R}$ » … \ill{} considérons alors …

Si
\marginal{vérifier s'il n'y a pas d'analogue des pts à val.\ dans une cat.\ 
réelle ???}
\note{la note de marge est oblique, en face du paragraphe « NB » ; le « Si »
isolé reste en suspens, suivi d'un double trait horizontal}

Pt \emph{motivique} de $X$ à valeurs dans un corps (ou un anneau) $k$ : les
objets de $T(k)$. L'analogue de la \emph{fibre d'un motif} \emph{en un
point} \ill{} les pts forment une catégorie fibrée sur (Schémas), munie de
fpqc, c'est un \uncertain{champ}. Pour $k$ donné, les isom. forment un
\emph{torseur}. Pour $k$ donné, les isom.\ d'un \struck{\ill{}} $\xi$ de
$T(k)$ sur un autre forment \uncertain{le} \struck{\ill{}}
\note{« Touseur » est souligné et biffé en tête de ligne, puis la phrase est
reprise ; elle se poursuit au feuillet 208, le 207 étant une page
dactylographiée étrangère}

\page{208}
\ill{} de \emph{classes de \ill{}}, ils sont [s'il y en a au moins un]
\uncertain{d'isomorphismes} \ill{} automorphismes des premiers, qui forment
les pts d'un groupe \uncertain{algébrique}. Mais en général il
\uncertain{peut} y avoir \ill{} des pts à valeurs dans un corps $k$ non
isomorphes. Un groupe de Galois opère \struck{\ill{}} sur l'ens.\ des pts à
valeurs dans $k$. Si $k$ est \ill{}, les pts à val.\ dans $k$, formant
\ill{}, \ill{} les \ill{} divisent \ill{} des \ill{} \ill{} $2$ (plus)
groupes \ill{} \emph{Aut}$(\xi)$.

\begin{itemize}
  \item Pts $\ell$-adiques : valeurs $\mathbb{Z}_\ell$ ou $\mathbb{Q}_\ell$
  \struck{\ill{}}
  \item Pts transcendants : valeurs $\mathbb{R}$ \struck{ou $\mathbb{C}$}
  voire $\mathbb{Z}$ [structure supplémentaire]. Il semble que ces pts
  correspondent biunivoquement aux pts de $X$ à valeurs dans $\mathbb{C}$,
  mod la conjugaison
  \item Pts Hodge-de Rham : valeurs dans $k$ quelconque
  \item Pts de Witt : valeurs $W_p(k)$ [$k$ corps de car.\ $p$]
\end{itemize}
\note{sur la page, les deux premières lignes sont réunies par une grande
accolade à gauche ; la dernière réserve est « corps de car.\ $p$ », ajoutée
au bout de la ligne des points de Witt}

\page{210}
\emph{Questions}
\begin{itemize}
  \item pts [de $X$] à val.\ dans $k$ [du cas $p > 0$] : correspondant aux
  [pts] de $X$ à val.\ dans $W(k)$ (ou corps des fractions ?), pts
  motiviques avec structure supplémentaire
  $\boxed{F, V}$ [foncteurs des groupes formels]
  \marginal{$\rightsquigarrow$ pts motiviques généralisés (sur une
  \uncertain{cat.\ divisée}) dans $k$, avec structure filtrée !}
  \item pts [de $X$] à val.\ dans $k$ de car.\ $0$ : correspondant aux pts
  motiviques dans $k$, avec \emph{structure filtrée}
  \item pts [de $X$] à valeurs en $\mathbb{C}$ mod la conj.\ complexe (=
  \uncertain{plus}) $\longleftrightarrow$ pts [motiviques] à valeurs dans
  $\boxed{\mathbb{R}}$, ou à valeurs dans $\boxed{\mathbb{C} \text{ avec } f_\infty}$,
  ou à valeurs dans $\mathbb{Z}$
\end{itemize}
\marginal{À voir}
\note{« $F, V$ », « $\mathbb{R}$ », « $\mathbb{C}$ avec $f_\infty$ » sont
entourés sur la page ; « F » a été écrit sur un premier signe. Une accolade
à gauche réunit les trois questions, et un long trait vertical descend dans
la marge jusqu'au bas du feuillet ; « À voir » est écrit obliquement dans la
marge gauche. La flèche ondulée part des points de car.\ $p > 0$}

\struck{\ill{}} \emph{cas $p$}

\emph{Pts banaux} : Pts à valeurs dans $\boxed{\mathbb{Z}_\ell}$, $\ell$
premier aux car.\ résiduelles, ou $\boxed{\mathbb{Q}_\ell}$.

\emph{Cas particulier important} : Pts à valeurs dans les corps finis
$\mathbb{F}_q$ correspondant aux pts à valeurs dans
$\boxed{W(\mathbb{F}_q)}$, avec structure supplémentaire $(F, V)$.
\note{les mots « Cas particulier important » et la fin de la ligne sont
soulignés}

\section{Groupe fondamental motivique discret}
\note{titre de sa main, souligné, en tête du feuillet 212}

\page{212}
Soit $\xi \in T(\mathbb{Z})$, on pose
\[
  \pi_1(X, \xi) = \mathrm{Aut}(\xi) = G_\xi(\mathbb{Z}) .
\]

\emph{N.B.} En vertu de Borel [\uncertain{-Harish}], il n'y a qu'un nombre
fini de $\xi \in T_\alpha(\mathbb{Z})$ possibles, pour $\alpha$ donné. Donc on
a un nb fini de groupes \emph{fondamentaux} \emph{discrets}. Sont-ils tous
\uncertain{isomorphes} \ill{} \struck{\ill{}} \ill{} ??

\[
  \xi \in T_{(\alpha)}(k), \quad k \text{ un anneau}
\]
$M$ motif $\in \mathrm{Ob}\,\mathcal{M}_\alpha$ ;

$M_\xi$ fibre en $\xi$, \quad $M_\xi = \xi(M)$ ;
\[
  \pi_1(X,\xi) = \mathrm{Aut}(\xi), \qquad G(X,\xi) =
  \underline{\mathrm{Aut}}(\xi) ;
\]
\struck{$G_{\cdot}$}
\[
  \pi_1(X,\xi) \simeq G(X,\xi)(k) .
\]
\note{les deux $\pi_1(X,\xi)$ de ces formules portent un indice peu lisible,
peut-être « $\mathrm{an}$ » ou « $0$ »}

Comparaison avec $\pi_1(X^{\mathrm{an}}, \xi)$ quand $\xi$ \uncertain{provient
d'un} \emph{pt} \ill{} …

\emph{Ex.} $X$ connexe sur $\mathbb{C}$, $\xi$ défini par un pt de $X$ à
valeurs dans $\mathbb{C}$. Le foncteur $\xi$ \struck{\ill{}} a une
structure \emph{filtrée}. $\xi$ considéré comme foncteur à valeurs dans les
modules \ill{} \ill{} $\pi_1(X,\xi)$ muni d'une structure de
\emph{Hodge} \ill{} \ill{} …

\page{214}
(Résulte de la conjonction des conj.\ de Hodge et de Tate [(sous forme
faible)], au moins si $X$ \uncertain{réduit}). \uncertain{Caractérisation}
\struck{de} de l'image ??
\note{« Caractérisation de l'image » est souligné deux fois ; un groupe de
cinq traits verticaux est tracé dans la marge gauche à la hauteur de ces
lignes}

Même conclusion si $X$ défini sur $k \subset \mathbb{C}$, $k$ alg.\ clos,
p.\ ex.\ $k = \overline{\mathbb{Q}}$.

\emph{Analyser l'exemple de Serre de $\pi_1$ non isomorphes}

[\emph{Cas corps de base $\mathbb{R}$} : tout pt \struck{\ill{}} [réel]
définit un $\xi_a \in T_\alpha(\mathbb{Z})$. \ill{} un $\xi$
\uncertain{isomorphe} [réel (\ill{} \ill{} $\xi$ \ill{} réel)]
\struck{\ill{} $X$ \ill{} \ill{} \ill{}}
\struck{\ill{} \ill{} \ill{} isom.}
Si \ill{} est un pt réel, alors $\xi$ \ill{} un \ill{} \uncertain{involution}
d'ordre $2$, qui \ill{} une \uncertain{autre} \ill{} de $G(X,\xi)$
\uncertain{compatible} avec la conj.\ \uncertain{complexe} de $G(X,\xi)$. Le
foncteur
\[
  M \mapsto \bigl(\xi(M_{\mathbb{C}}),\ \text{structure de Hodge} +
  \text{réalification définie par } f_\infty\bigr)
\]
est pleinement fidèle.]

\begin{tikzcd}[column sep=small, row sep=small]
 & X \arrow[dd, no head] \\
\mathbb{C} \arrow[ur, no head] \arrow[dr, no head] & \\
 & \mathbb{R}
\end{tikzcd}

\note{petit schéma dans la marge gauche, à la hauteur du « Cas corps de base
$\mathbb{R}$ » : $X$ au-dessus de $\mathbb{R}$, $\mathbb{C}$ à gauche, trois
traits sans flèche}

\emph{Pb} étudier comment varie la structure de Hodge pour un pt [complexe]
variable, \emph{en fonction de la donnée} \uncertain{définie} en le pt
\uncertain{$\infty$},
\note{un groupe de trois traits verticaux marque ces lignes dans la marge ;
la phrase s'arrête sur la virgule au bas du feuillet}

\page{216}
Lien avec la variété des modules pour des motifs polarisés rigidifiés : si
$\mathcal{M}$ est une telle variété modulaire, alors un morphisme $X \to
\mathcal{M}$ est déterminé quand on connaît sa valeur en un seul point,
plus l'application $\pi_1(X) \to \pi_1(\mathcal{M})$ …

C'est une propriété intrinsèque importante de $\mathcal{M}$, et une
propriété de cette nature mérite d'être étudiée pour les groupes
\uncertain{unités} … Peut-être \uncertain{les} \uncertain{\ill{}}
\struck{\ill{}} [intervention] \ill{} dites \struck{\ill{}} [variétés]
modulaires pour certains types de motifs.

Introduire la notion de schéma \uncertain{pointé} [motiviquement], [(i.e.\ 
dans un anneau $k$)] et morphismes de tels, qui définissent des morphismes
correspondants pour les
\[
  G(X,\xi) \to G(Y,\eta),
\]
donc pour
\[
  G(X,\xi)(k) = \pi_1(X,\xi) \to G(Y,\eta)(k) = \pi_1(Y,\eta) .
\]

\page{218}
\note{feuillet très repris : lignes biffées, insertions reliées par des
boucles ; on transcrit la couche qui subsiste}

$X$ \struck{\ill{}} \uncertain{variété} de type fini
[\uncertain{univers.}] \uncertain{universellement} \ill{} sur un corps $k$,
$v, w$ deux places archimédiennes de \struck{$k$} $X$, d'où foncteurs fibres
$\xi_v^{\mathbb{Z}}, \xi_w^{\mathbb{Z}}$. Question : a) [Sont-ils]
isomorphes ? b) \struck{N'est-il \ill{} que} [\uncertain{Quand} \ill{}
\ill{}] [mult.] \ill{} pour les sous-catégories de t.f.\ de
$\mathcal{M}(X)$, les restrictions de ces foncteurs à
$\mathcal{M}_\alpha(X)$ soient isomorphes ?

N.B. Nous pouvons supposer, quitte à remplacer $k$ par
\struck{\uncertain{sa clôture}} la clôture alg.\ de $k$ dans $k(X)$, que $k$
est alg.\ clos dans $k(X)$ [utiliser le fait $X$ \uncertain{universel} !],
donc $X$ géom.\ irréd.\ sur $k$. Soit $k_0$ la clôture alg.\ de
$\mathbb{Q}$ dans $k$.
\note{suivent deux lignes et demie biffées et encadrées d'une boucle, où se
lisent « le corps », « des places (= \ill{}) », « \ill{} sur $k_0$ » ; puis,
au-dessus d'une ligne « 1) $k_0$ \ill{} fini dans $k$ » également biffée,
l'insertion « contient la clôture alg.\ $E$ de $\mathbb{Q}$ dans $k$
(p.\ ex.\ $k_0$ alg.\ fini dans $k$). Alors »}

1°) \uncertain{si $X$ est défini sur $k_0$.}

a) $\xi_v^{\mathbb{Z}} \mid \mathcal{M}_\alpha(X) \simeq
\xi_w^{\mathbb{Z}} \mid \mathcal{M}_\alpha(X)$ pour $\forall \alpha$.

En effet, $X$ et $\mathcal{M}_\alpha(X)$ proviennent de $X_0$,
[$\mathcal{M}_\alpha(X_0)$], de type fini \ill{}
\struck{$X_0 \to \operatorname{Spec} A_0$ fid.\ plat \ill{}}
[$A_0 \subset k$] de type fini sur $k_0$, \struck{et} $\mathbb{Q}$ est
canoniquement \ill{} ; on en déduit \ill{} $k = E_0$, \ill{}
\ill{} on obtient le résultat voulu grâce au groupe fondamental
transcendant. \struck{\ill{}}
\note{« $\mathcal{M}_\alpha(X_0)$ » est au-dessus de la ligne ; « $A_0$ »,
« $k_0$ » sont écrits en surcharge d'autres lettres}

\begin{tikzcd}[column sep=small, row sep=small]
 & X_0 \arrow[r, no head] \arrow[d, no head] & X \arrow[d, no head] \\
k_0 \arrow[r, no head] & A_0 \arrow[r, no head] & k
\end{tikzcd}

\note{schéma de la marge gauche ; au-dessous, un second schéma du même
genre ($X_0$, $X$, $E$, $k_0$, $A_0$, $k$, $\mathbb{Q}$) est couvert de
hachures et de traits, illisible}

b) Si \struck{\ill{}} [$v, w$ coïncident sur $k_0$], \ill{} \ill{}
$\xi_v^{\mathbb{Z}} \simeq \xi_w^{\mathbb{Z}}$.
\struck{sinon, \ill{} \ill{} $\xi_v^{\mathbb{Z}} \not\simeq
\xi_w^{\mathbb{Z}}$ ; \ill{}}
\note{deux rédactions superposées de b), la seconde biffée ; on transcrit
ce qui survit, avec réserve}

\marginal{on peut supposer $X_0 \to \operatorname{Spec} A_0$ fid.\ plat
\uncertain{à fibres connexes}, et $\operatorname{Spec}(A_0) \to k_0$ à
fibres connexes, donc $X_0 \to \operatorname{Spec} k_0$ à fibres connexes.}
\note{note en bas du feuillet, marquée d'un trait à gauche}

\page{220}
[\uncertain{dénombrable} (p.\ ex.\ de t.f.\ sur $\mathbb{Q}$)]
\struck{\ill{} \ill{} de génération \ill{}}, on voit encore que
$\xi_v^{\mathbb{Z}} \simeq \xi_w^{\mathbb{Z}}$. En effet, $k$ est réunion
filtrante d'une \struck{suite} [\ill{}] $A_i$ de sous-$k_0$-algèbres de type
fini, où l'on peut supposer les morphismes de transition
$\operatorname{Spec} A_j \to \operatorname{Spec} A_i$ lisses à fibres géom.\ 
connexes, donc aussi les $X_j \to X_i$, donc aussi les $X_{j\mathbb{C}} \to
X_{i\mathbb{C}}$. Or si des pts $x_i, y_i$ des $X_{i\mathbb{C}}$,
s'envoyant les uns sur les autres par les morphismes de transition, et
qui \ill{} la limite \ill{} de la connexité des fibres, \ill{} pour les
morphismes de transition
\[
  \pi^1(x_j, y_j) \to \pi^1(x_i, y_i)
\]
sont surjectifs. Donc $\varprojlim \pi^1(x_i, y_i) \neq \emptyset$, cqfd.

Si $k$ pas dénombrable, [\uncertain{on se ramène au cas}] \ill{} sur
$E$.

2) $E$ \struck{$\neq k_0$ alg.\ fini dans $k$}. Alors je
\uncertain{pense} qu'il existe $\alpha$ tel que \struck{\ill{}}
\[
  \xi_v^{\mathbb{Z}} \mid \mathcal{M}_\alpha(X) \not\simeq
  \xi_w^{\mathbb{Z}} \mid \mathcal{M}_\alpha(X) .
\]
On peut supposer que $X = \operatorname{Spec}$ \uncertain{$k$}, $k$ extension
finie de $\mathbb{Q}$.
\note{la lettre après « Spec » est écrite en surcharge d'un $E$. Le reste du
feuillet --- trois lignes et un cadre --- est biffé ; on y lit
« $\mathbb{Q}$ » à deux reprises. Dans la marge gauche, une formule en
$\xi_v^{\mathbb{Z}}(M)$, $\xi_w^{\mathbb{Z}}(M)$ est couverte de traits.
Le lot 12 s'ouvre (feuillet 222) sur « Il s'agirait de construire $X$
projectif lisse sur $k$ … », qui en est vraisemblablement la suite}

\end{document}
